\documentclass[11pt]{scrartcl}

% =============================================================================
% PACKAGES
% =============================================================================

\usepackage[english]{babel}
\usepackage[T1]{fontenc}
\usepackage{fourier}
\usepackage{geometry}
\usepackage{enumitem}
\usepackage{scrlayer-scrpage}
\usepackage{xcolor}

% =============================================================================
% PAGE LAYOUT AND TYPOGRAPHY
% =============================================================================

\geometry{
  paper=a4paper,
  top=2.5cm,
  bottom=3cm,
  left=3cm,
  right=3cm,
  headheight=0.75cm,
  footskip=1.5cm,
  headsep=0.75cm
}

% Paragraph formatting
\setlength{\parindent}{0pt}
\setlength{\parskip}{4pt}

\setlist[itemize]{noitemsep}

\setkomafont{section}{\normalfont\scshape}
\setkomafont{subsection}{\normalfont\bfseries}
\setkomafont{subsubsection}{\normalfont\itshape}
\setkomafont{paragraph}{\normalfont\scshape}

\clearpairofpagestyles
\cfoot*{\pagemark}

\pagestyle{scrheadings}


% =============================================================================
% ENTRY COMMANDS
% =============================================================================

% Publications
\newcommand{\pubentry}[5]{%
  \item
  #1.
  \textbf{#2}.
  #3.
  #4.

  \begin{sloppypar}
    \vspace{8pt}
    \textit{#5}
    \vspace{8pt}
  \end{sloppypar}
}

% Conference / International Meeting
\newcommand{\confentry}[4]{%
  \item
  \textbf{#1} -- #2, #3, #4.
}

% Course / School
\newcommand{\courseentry}[7]{%
  \item
  \textbf{#1} by #2, #3.\\
  Period: #4. Duration: #5. Credits: #6 CD.\\
  Exam type: #7.
}

% Tutoring
\newcommand{\tutorentry}[6]{%
  \item
  Tutor for the \textbf{#1} course (cod. #2, #3).\\
  Course: #4.\\
  Academic year: #5. Duration: #6.
}

% Period Abroad
\newcommand{\abroadentry}[5]{%
  \item
  \textbf{#1}.\\
  Host institution: #2.\\
  Period: #3.\\
  Abroad supervisor: #4.\\
  Activity: #5.
}

% Conferences organization
\newcommand{\conforgentry}[4]{%
  \item
  \textbf{#1} -- #2, #3, #4.
}


% =============================================================================
% SECTION ENVIRONMENTS
% =============================================================================

% Publications
\newenvironment{publications}
  {%
    \subsection*{Publications}
    \begin{itemize}
  }
  {%
    \end{itemize}
  }

% Conferences and International Meetings
\newenvironment{conferences}
  {%
    \subsection*{Conferences and International Meetings}
    \begin{itemize}
  }
  {%
    \end{itemize}
  }

% Courses and Schools
\newenvironment{courses}
  {%
    \subsection*{Courses and Schools}
    \begin{itemize}
  }
  {%
    \end{itemize}
  }

% Tutoring Activities
\newenvironment{tutoring}
  {%
    \subsection*{Tutoring Activities}
    \begin{itemize}
  }
  {%
    \end{itemize}
  }

% Conferences organization
\newenvironment{conforg}
{%
  \subsection*{Conferences Organization}
  \begin{itemize}
}
{%
  \end{itemize}
}

% Period Abroad
\newenvironment{periodabroad}
{%
  \subsection*{Period Abroad}
  \begin{itemize}
}
{%
  \end{itemize}
}



% =============================================================================
% DOCUMENT INFORMATION
% =============================================================================

\title{%
  \normalfont\normalsize

  \textsc{%
    Università di Bologna,
    Dipartimento di Informatica -- Scienza e Ingegneria
  }\\

  \vspace{25pt}
  \rule{\linewidth}{0.5pt}\\
  \vspace{20pt}

  {\huge PhD Annual Report -- Second Year\\
  \vspace{15pt}
  \LARGE Advances in Collective Robotics Through Macro-Programming\\[14pt]
  }

  \vspace{12pt}
  \rule{\linewidth}{2pt}\\
  \vspace{12pt}
}

\author{%
  \LARGE Angela Cortecchia\\[14pt]
  \normalsize
  \begin{tabular}{@{}ll@{}}
    Supervisor: & Prof. Danilo Pianini \\
    Co-supervisor: & Prof. Mirko Viroli \\
    Committee member: & Prof. Enrico Gallinucci
  \end{tabular}%
}

\date{}


% =============================================================================
% DOCUMENT
% =============================================================================

\begin{document}

\maketitle

\subsection*{Research}

During the second year of my PhD, my research focused on advances in collective robotics through macro-programming,
with Aggregate Computing (AC) as the main programming paradigm adopted throughout the work.
%
Collective robotic systems consist of multiple interacting devices that must coordinate their local actions to achieve system-level objectives,
often while operating under mobility,
failures, and continuously changing communication conditions.
%
In these settings, relying on fixed roles or predefined coordination structures can limit the ability of the system to adapt at runtime.

Aggregate Computing provides a different perspective by expressing collective behavior through distributed computations that emerge from local interactions among devices.
%
My research investigates how this programming model can be exploited to engineer robotic collectives whose organization can adapt together with the system and its environment.
%
During this year,
I studied this problem from complementary perspectives,
ranging from resilient information propagation to distributed state estimation, morphogenetic organization, and runtime safety.
%
Although addressing different aspects of collective robotics, these activities share the objective of developing reusable mechanisms for the coordination and management of collective robotic systems.

\paragraph*{Self-stabilizing Gossip and Resilient Coordination}

I completed a self-stabilizing min--max gossip algorithm for Aggregate Computing, designed to recover from transient faults and topology changes without explicit resets.
%
The approach relies on path--loop detection to remove obsolete information and avoid cyclic dependencies,
while preserving lightweight local communication.
%
The algorithm was formally proved self-stabilizing,
implemented in \textit{Collektive}, and published at COORDINATION 2026.

The work is currently being extended to generalize the underlying mechanism beyond min--max consensus and investigate its applicability to a broader range of resilient collective coordination problems.

\paragraph*{Self-organizing Distributed Particle Filtering}

Another major research line concerned Distributed Particle Filtering (DPF).
%
Existing DPF approaches distribute the estimation process according to different predefined organizations.
%
For instance, some rely on a central fusion node that collects and combines information from the other sensors,
while others assign coordination responsibilities to selected leaders or use fully decentralized consensus-based schemes.
%
These approaches provide different trade-offs, but their coordination structure, node roles, and information flows are generally defined as part of the filtering protocol itself.
%
My work investigates instead whether these aspects can be expressed independently of the filtering logic and adapted at runtime as the sensing system changes.

To this end, the main components of DPF are expressed through computational fields,
namely distributed data structures that associate values with the devices of the system and evolve over space and time through local interactions.
%
This representation makes it possible to express both estimation and coordination at the collective level,
allowing the organization of the filtering process to change when connectivity, sensing conditions, mobility, or device availability change.

The first formulation,
presented at DCOSS-IoT 2026,
focused on showing how different DPF organizations could be expressed within the same field-based model by changing the coordination and information-dissemination strategies rather than the underlying filtering logic.
%
The work was subsequently extended for ACSOS 2026 towards a self-organizing multi-target tracking scenario.
%
In this setting, multiple moving targets, represented by zebras in the simulated scenario,
are observed by a distributed sensing system that also includes mobile observers.
%
The mobile observers adapt their spatial organization as the targets move, while the DPF continuously maintains their estimates.
%
This allowed us to study distributed estimation together with mobility and changing network connectivity,
showing how sensing and coordination roles can adapt at runtime within the same field-based formulation.

\paragraph*{Safety Guarantees for Aggregate-Programmed Robot Swarms}

From January to April 2026, I carried out a research period at Mälardalen University (MDU),
within the Mälardalen University Automation Research Center (MARC),
under the supervision of Prof. Alessandro Vittorio Papadopoulos.

The activity focused on integrating Control Lyapunov Functions (CLFs) and Control Barrier Functions (CBFs) with Aggregate Computing.
%
The resulting architecture separates the nominal collective behavior, computed by the aggregate program,
from a distributed safety-filtering layer that modifies the resulting commands whenever they would violate safety requirements.
%
This allows the swarm to pursue its collective objective while accounting for constraints related to obstacles,
nearby robots, and communication connectivity.
%
A working implementation was developed in \textit{Collektive} and validated through simulation in Alchemist.

The collaboration with the MARC group continued after the research stay and led to a joint paper presented at SISSY 2026.
%
Further work is planned on more complex formations, such as geometric formations that must be preserved while the swarm moves collectively through the environment,
and on the corresponding safety guarantees during their adaptation.

\paragraph*{Morphogenetic Self-Organization}

The research on FieldVMC also continued during the year and was consolidated with the publication of the journal article in
\textit{Complex \& Intelligent Systems}.
%
FieldVMC provides a decentralized and asynchronous formulation of morphogenetic growth over arbitrary network topologies.

I am now extending FieldVMC by combining it with Artificial Potential Fields (APFs) and Signed Distance Fields (SDFs).
%
The idea is to use APFs to steer the local evolution of the structure, while SDFs provide a geometric description of the target shape and its boundaries.
%
Together, these mechanisms are being investigated as a way to guide morphogenetic growth towards more complex artificial structures in simulation without losing the decentralized and self-organizing properties of the original approach.

\paragraph*{Collektive Framework and Coordination Library}

In parallel with the development of these algorithms,
I am contributing to an ongoing journal work on \textit{Collektive},
the Kotlin-based framework used throughout my research to implement Aggregate Computing applications.
%
The work aims to consolidate Collektive as a modern framework for programming heterogeneous collective systems,
while providing a library of reusable coordination abstractions on top of its core programming model.

The mechanisms developed throughout my PhD are intended to progressively contribute to this library rather than remain isolated implementations tied to individual experiments.
%
In particular, the self-stabilizing coordination mechanisms and other reusable abstractions emerging from the work on distributed estimation,
morphogenesis, and collective robotics can become part of the higher-level functionality offered by Collektive.
%
This provides a direct path from the research results to reusable software components that can support the development of new aggregate applications.

\paragraph*{Scientific Community Service}

During 2026, I served as Publicity Chair for the 7th IEEE International Conference on Autonomic Computing and Self-Organizing Systems (ACSOS 2026), hosted at the Cesena Campus of the University of Bologna. I also contributed to the scientific community as a reviewer for \textit{Complex \& Intelligent Systems}, a Q1 peer-reviewed journal.

\paragraph*{Research Plan for the Next Year}

During the next year,
I plan to continue the main research directions developed so far while making them more tightly connected.
%
The FieldVMC work will be extended with APF- and SDF-based guidance to investigate the generation of more complex artificial structures.
%
The collaboration with Mälardalen University will also continue, extending the current safety-filter architecture towards more advanced swarm formations and safety constraints.

The self-stabilizing gossip work will be further developed by generalizing the mechanism introduced for min--max consensus and investigating its use for other collective coordination problems.

The work on self-organizing DPF will be extended towards more realistic and heterogeneous distributed sensing scenarios.
%
In particular, I plan to investigate systems in which different types of simulated devices coexist,
such as static sensing infrastructure and mobile robotic sensors, and cooperate within the same estimation process.
%
Mobile sensors will not only contribute observations,
but will also coordinate their motion to maintain a stable spatial formation while adapting to the targets and to changes in the sensing environment.
%
This will allow the approach to be evaluated in more realistic simulation scenarios,
where distributed estimation, heterogeneous sensing, mobility, and collective motion must operate together rather than being studied in isolation.

In parallel,
I will continue contributing to the development of \textit{Collektive} and to the ongoing journal work on the framework.
%
A particular objective will be to consolidate the reusable mechanisms developed throughout the PhD into the coordination libraries offered by Collektive,
making them available as composable building blocks for the development of aggregate applications.

A further objective will be to investigate how multiple collective applications can coexist and be adapted or preempted at runtime over partially overlapping and dynamically evolving groups of devices.
%
This moves the focus from individual collective algorithms towards the runtime management and interaction of multiple collective processes.

% =============================================================================
% COURSES AND SCHOOLS
%
% \courseentry
% {title}
% {professor}
% {location}
% {period}
% {duration}
% {credits}
% {exam type}

% =============================================================================

\begin{courses}

  \courseentry
  {Literature Reviews: Methods and Tools}
  {Prof. Roberto Casadei}
  {Department of Computer Science and Engineering, University of Bologna}
  {April -- May 2026}
  {12 hours}
  {2.4}
  {Project (passed)}

  \courseentry
  {Foundations and Emerging Trends in Serverless Computing}
  {Prof. Andrea Sabbioni}
  {Department of Computer Science and Engineering, University of Bologna}
  {January -- February 2026}
  {10 hours}
  {0.4}
  {Attendance}

  \courseentry
  {Data Platforms and Artificial Intelligence: Challengesd and Applications}
  {Prof. Matteo Francia}
  {Department of Computer Science and Engineering, University of Bologna}
  {November 2025}
  {12 hours}
  {0.48}
  {Attendance}

\end{courses}

% =============================================================================
% CONFERENCES and OTHER INTERNATIONAL MEETINGS
%
% \confentry
% {acronym}
% {title}
% {date}
% {location}
% =============================================================================

\begin{conferences}

  \confentry
  {COORDINATION 2026}
  {21st International Conference on Coordination Models and Languages}
  {9th -- 11th June 2026}
  {Urbino, Italy}

  \confentry
  {WOA 2026}
  {27th Workshop From Objects to Agents}
  {15th -- 17th June 2026}
  {Salerno, Italy}

  \confentry
  {DCOSS-IoT 2026}
  {22nd Annual International Conference on Distributed Computing in Smart Systems and the Internet of Things}
  {22nd -- 24th June 2026}
  {Reykjavík, Iceland}

\end{conferences}

% =============================================================================
% Conferences organization

\begin{conforg}
  \confentry{ACSOS 2026 Publicity Chair}{7th IEEE International Conference on Autonomous Computing and Self-Organizing Systems}{7th -- 11th September, 2026}{Cesena, Italy}
\end{conforg}



% =============================================================================
% TUTORING ACTIVITIES
%
% \tutorentry
% {title}
% {code}
% {professor}
% {course}
% {academic year}
% {duration}
% =============================================================================

\begin{tutoring}

  \tutorentry
  {Object Oriented Programming}
  {\textbf{70219}}
  {Prof. Mirko Viroli}
  {Bachelor's Degree Program}
  {2025--2026}
  {60 hours}

  \tutorentry
  {Object Oriented Programming}
  {\textbf{70219}}
  {Prof. Mirko Viroli}
  {Bachelor's Degree Program}
  {2026--2027}
  {60 hours}

\end{tutoring}

% =============================================================================
% PERIOD ABROAD
%
% \abroadentry
% {title}
% {host institution}
% {period}
% {abroad supervisor}
% {activity}
% =============================================================================

\begin{periodabroad}

  \abroadentry
  {Integration of Control Barrier Functions for Safety-Critical Guarantees in Aggregate Computing}
  {MARC -- Mälardalen University Automation Research Center, Mälardalen University (MDU), Sweden}
  {January -- April 2026}
  {Prof. Alessandro Vittorio Papadopoulos}
  {Research on the integration of Control Lyapunov Functions and Control Barrier Functions into Aggregate Computing, with a focus on providing stability and safety guarantees during transient adaptation phases. The activity included the design and implementation of a working integration within the Aggregate Computing framework, the investigation of mappings from collective-level safety requirements to local control constraints, and the simulation-based validation of safety-critical multi-robot scenarios.}

\end{periodabroad}

% =============================================================================
% PUBLICATIONS
%
% \pubentry
% {authors}
% {title}
% {venue}
% {date}
% {abstract}
% =============================================================================

\begin{publications}

  \pubentry
  {\textbf{Angela Cortecchia}, Danilo Pianini, and Mirko Viroli}
  {A Self-stabilizing Min--Max Consensus via Path--Loop Detection}
  {International Conference on Coordination Models and Languages (COORDINATION 2026)}
  {June 2026}
  {In large-scale distributed systems, such as the Internet of Things (IoT), min-max consensus algorithms provide a mechanism for collective coordination by enabling nodes to converge on a “best” value produced by one of the participants in the computation. However, min-max consensus algorithms are monotonic and non-self-stabilizing by nature: once a value is merged into the aggregate it cannot be retracted, leading to propagation of stale or incorrect data in the presence of transient faults or topology changes. In this work, we propose a novel self-stabilizing min-max consensus algorithm ensuring convergence to the best available value in the network by propagating information along shortest valid paths. Each gossip message carries a value and path of nodes that have acknowledged it, enabling loop-freedom and natural pruning of obsolete contributions. We rely on field-based coordination and specifically the Aggregate Computing paradigm to present the algorithm, prove self-stabilization, and provide an implementation as a reusable library for the \textit{Collektive} DSL. This work contributes a foundational building block for resilient coordination in pervasive computing systems, paving the way to more complex, self-stabilizing distributed applications.}

%  \pubentry
%  {\textbf{Angela Cortecchia}, Giovanni Ciatto, Roberto Casadei, and Danilo Pianini}
%  {FieldVMC: an asynchronous model and platform for self-organising morphogenesis of artificial structures}
%  {\emph{Complex \& Intelligent Systems}, vol.~12, no.~2}
%  {2026}
%  {Introduces FieldVMC, an asynchronous model and platform for the self-organising morphogenesis of artificial structures.}

  \pubentry
  {\textbf{Angela Cortecchia}, Davide Domini, Giovanni Ciatto, Roberto Casadei, Danilo Pianini, and Mirko Viroli}
  {Flexible Distributed Particle Filtering for the Internet of Things via Aggregate Computing}
  {22nd International Conference on Distributed Computing in Smart Systems and the Internet of Things (DCOSS-IoT 2026)}
  {June 2026}
  {State estimation from uncertain, distributed observations is central in many cyber-physical applications. While Distributed Particle Filtering (DPF) algorithms address nonlinear and non-Gaussian estimations in distributed settings, most solutions remain tied to specific architectures and communication assumptions, limiting adaptability in open, heterogeneous deployments-most notably, the Internet of Things (IoT). In this paper, we propose a field-based formulation of Distributed Particle Filtering grounded in Aggregate Computing (AC). By expressing estimation and information dissemination as computational fields, our approach decouples the core filtering logic from coordination and data-flow strategies. This enables systematic customisation of key design dimensions, including fusion-center placement and resilience, aggregated measurement functions, as well as the type and scope of information propagation. Through a set of in-silico experiments, we show how diverse DPF configurations can be derived within a unified framework, highlighting trade-offs among accuracy, communication cost, and robustness. Overall, the proposed approach positions AC as an effective abstraction layer for engineering adaptable DPF solutions in open IoT environments.}

  \pubentry
  {Manuel Andruccioli, \textbf{Angela Cortecchia}, Davide Domini, Nicolas Farabegoli, Giovanni Delnevo, Danilo Pianini, Riccardo Venanzi, and Mirko Viroli}
  {HarmoniKt: a Unifying Middleware for Heterogeneous Robot Fleets}
  {2026 IEEE 23rd Consumer Communications \& Networking Conference (CCNC)}
  {January 2026}
  {In recent years, companies have increasingly invested in Industry 4.0 research to automate repetitive tasks or activities that may be harmful to humans. For instance, large-scale warehouses, such as those operated by Amazon, rely heavily on mobile robots to streamline logistics operations and ensure worker safety. At the same time, robot manufacturers are releasing more reliable platforms with advanced capabilities, making them suitable for complex industrial tasks. Despite this growing interest and technological progress, significant challenges remain. A key and non-trivial issue lies in the integration of heterogeneous robot fleets: each vendor typically employs proprietary technologies and interfaces, which hinders interoperability and limits the potential of multi-brand deployments.In this work, we introduce HarmoniKt, an extensible middle-ware that addresses this challenge by introducing an abstraction layer and providing a unified REST API for the control and management of heterogeneous robots. Our solution has been validated in a physical industrial-like environment using a mixed fleet of Boston Dynamics Spot and Mobile Industrial Robots (MiR). Furthermore, we present a comparative analysis showing that the access latency introduced by our middleware is not significantly higher than that of direct robot access, demonstrating the feasibility of unified robot management without compromising performance.}

  \pubentry
  {Filippo Gurioli, Martina Baiardi, \textbf{Angela Cortecchia}, and Danilo Pianini}
  {High-Fidelity Simulation of Aggregate Computing Systems with Collektivity}
  {International Conference on Coordination Models and Languages (COORDINATION 2026)}
  {June 2026}
  {The verification and validation of collective adaptive systems is a challenging task, typically tackled through a combination of approaches including formal analysis, simulation, Hardware-in-the-Loop, and real-world experimentation. One critical aspect of simulation concerns the trade-off between realism and scalability. A typical validation pipeline involves initial simulation of large-scale scenarios with simplified models of the world, and detail is gradually increased in subsequent steps, introducing more realistic physical and environmental dynamics. An underexplored opportunity in this context lies in leveraging game development platforms to achieve high-fidelity simulations. In this work, we explore the integration of an existing implementation of aggregate programming, a common paradigm for engineering collective adaptive systems, with Unity, a widely used game development platform. We show that this integration is feasible even though the two systems were not designed to work together, we discuss the technical challenges that were encountered and the limitations of the approach, while highlighting the potential for such integration to enhance simulation realism in the study of collective adaptive systems.}

  \pubentry
  {\textbf{Angela Cortecchia}, Davide Domini, Giovanni Ciatto, Roberto Casadei, and Mirko Viroli}
  {Multi-Target Tracking via Field-Based Distributed Particle Filtering}
  {7th IEEE International Conference on Autonomic Computing and Self-Organizing Systems (ACSOS 2026)}
  {September 2026}
  {Distributed Particle Filtering (DPF) is effective for nonlinear and non-Gaussian state estimation in distributed sensing systems, but existing approaches are often tied to fixed architectures, predefined roles, or rigid communication assumptions. This limits their applicability in open and dynamic scenarios, where sensors may be heterogeneous, mobile, intermittently connected, or subject to failures. In this paper, we reinterpret DPF through Aggregate Computing (AC), modelling sensing, information dissemination, role assignment, and particle-based estimation as computational fields. This decouples filtering logic from coordination mechanisms, turning architectural choices into composable and runtime-adaptive aggregate patterns. The resulting formulation supports self-organising DPF: nodes interact locally while collectively maintaining estimates, adapting sensing regions, and dynamically assigning roles such as leaders or regional coordinators. It also enables self-healing under failures and topology changes, and supports interoperability among fixed and mobile sensing nodes. We evaluate the approach in a multitarget tracking scenario with mobile observers, showing that field-based DPF supports robust and flexible estimation under changing connectivity and spatial configurations. Overall, AC emerges as a unifying abstraction layer for engineering adaptive and self-organising DPF solutions in open, heterogeneous, and dynamic Internet of Things (IoT) environments.}

  \pubentry
  {\textbf{Angela Cortecchia}, Alessandro Papadopoulos, and Danilo Pianini}
  {Toward Safe Aggregate Computing: A Distributed Control-Theoretic Safety Filter for Robot Swarms}
  {7th IEEE International Conference on Autonomic Computing and Self-Organizing Systems (ACSOS 2026)}
  {September 2026}
  {Robot swarms must coordinate their actions to achieve global objectives in applications such as search and rescue or environmental monitoring. In these scenarios, unsafe transient behavior may cause collisions with other robots or environmental obstacles or loss of connectivity before the swarm reaches its intended configuration. This paper proposes an architecture that integrates Aggregate Computing with a control-theoretic safety filter for robot swarms. The aggregate program computes the nominal collective behavior, while a safety-filtering layer refines the resulting commands through Control Lyapunov Function and Control Barrier Function constraints. This layer translates convergence and safety requirements into constraints that are enforced on each robot in a distributed fashion. We evaluate the approach in proof-of-concept simulations involving multiple targets, leader-based coordination, connectivity preservation, and environments with static obstacles. The results illustrate how Aggregate Computing can specify adaptive swarm-level objectives, while the safety filter modifies unsafe nominal commands to respect constraints whenever the active constraints are feasible.}

  \pubentry
  {\textbf{Angela Cortecchia}}
  {Towards Collective Robotic Operating Systems through Aggregate Computing}
  {7th IEEE International Conference on Autonomic Computing and Self-Organizing Systems (ACSOS 2026)}
  {September 2026}
  {Swarm robotic systems are increasingly relevant for missions in which many heterogeneous devices must operate in dynamic and potentially safety-critical environments. Programming such systems remains challenging: developers need to express objectives at the collective level, while their execution is distributed across devices that may move, fail, join or leave the system, without a centralized point of coordination. Aggregate Computing addresses this issue by providing macroprogramming abstractions for specifying resilient and self-organizing collective behaviors. However, current Aggregate Computing applications typically execute a single collective program and offer limited support for runtime adaptation, preemption, concurrent applications, and safety guarantees during transient phases. This research investigates the foundations of a Collective Robotic Operating System built on Aggregate Computing: a middleware layer that lifts ROS-inspired capabilities from individual robots to the collective level. The paper outlines the main requirements for such a system and reports ongoing contributions on morphogenetic formation control, runtime path replanning, self-stabilizing consensus, collective sensing, and safety guarantees. These results provide initial building blocks towards a collective operating layer for swarm robotics, with future work focusing on safe runtime adaptation and support for multiple concurrent collective processes.}

\end{publications}

\end{document}